\documentclass[11pt,a4paper]{ctexart}
\usepackage[margin=2.2cm]{geometry}
\usepackage{amsmath,amssymb,amsthm,mathtools,booktabs,hyperref,microtype}
\hypersetup{colorlinks=true,linkcolor=black,urlcolor=blue,citecolor=black}
\newtheorem{theorem}{定理}
\newtheorem{lemma}{引理}
\newtheorem{corollary}{推论}
\renewcommand{\proofname}{证明}
\newcommand{\F}{F_{3,D}}
\newcommand{\vthree}{v_3}
\title{\bfseries 多输入 \(F_{3,D}\) 的固定多项式实现\\与整数热带表达式}
\author{研究札记 · F3D-MULTI}
\date{2026 年 9 月 29 日}
\begin{document}
\maketitle

\begin{abstract}
在
\[
\mathcal D=\{(n,d)\in\mathbb Z^2:n\ne\pm d\}
\]
上令
\[
F_{3,D}(n,d)=v_3(n+d)-v_3(n-d).
\]
本文把固定整数多项式对的单输入表达能力分类推广到任意固定有限输入数。首先给出二次正性检测器，并构造双输入最小值/最大值以及批量最小赋值与第一极小位置的固定多项式实现。主结果证明：函数 \(g:\mathbb Z^m\to\mathbb Z\) 可由一个全域有效的固定整数多项式对实现，当且仅当 \(g\) 可由整数常数、坐标、加减、最小值与最大值有限次复合；等价地，\(g\) 是两个整数热带多项式之差。必要性证明使用实际整数输入上的有限单位网格与 Lagrange 插值，控制最低项抵消并将多项式赋值化为有限整数仿射分区函数。由分类进一步得到统一 Lipschitz 界，从而排除无界幅度的条件开关，而固定幅度开关仍可实现。
\end{abstract}

\section{和差比坐标}
对输入 \(X=(n,d)\in\mathcal D\)，令
\[
u=n+d,\qquad v=n-d,\qquad x=u/v.
\]
则
\[
\boxed{\F(X)=v_3(x).}
\]
任意非零有理数 \(x=a/b\) 都可由
\[
n=a+b,\qquad d=a-b
\]
实现。

\section{二次正性检测}
定义
\[
\Theta(n,d)=
\left(
3n^2-2nd+3d^2,\,
(n-d)^2
\right).
\]
在 \(u=n+d,v=n-d\) 坐标下，
\[
\Theta=(u^2+2v^2,v^2),
\]
所以
\[
N+D=u^2+3v^2,\qquad N-D=u^2+v^2.
\]
若
\[
a=v_3(u),\qquad b=v_3(v),
\]
则
\[
v_3(u^2+3v^2)=\min(2a,1+2b)
\]
因为两个候选赋值奇偶不同；同时
\[
v_3(u^2+v^2)=2\min(a,b).
\]
因此
\[
\boxed{
\F(\Theta(n,d))
=
\begin{cases}
1,&a>b,\\
0,&a\le b.
\end{cases}
}
\]
亦即
\[
\boxed{
\F(\Theta(n,d))=\mathbf1_{\{\F(n,d)>0\}}.
}
\]
并且 \(N>D>0\)。

\section{双输入最小值与最大值}
对非零有理数 \(x,y\)，定义
\[
\mu(x,y)=
\frac{x^3+3y^3}{x^2+3y^2}.
\]
分子和分母在 \((\mathbb Q^\times)^2\) 上均不为零。

令
\[
a=v_3(x),\qquad b=v_3(y).
\]
有
\[
v_3(x^3+3y^3)=\min(3a,1+3b),
\]
\[
v_3(x^2+3y^2)=\min(2a,1+2b),
\]
两边都没有并列最低项。

分情况即得
\[
\boxed{v_3(\mu(x,y))=\min(a,b).}
\]
因此
\[
\boxed{
v_3\!\left(\frac{xy}{\mu(x,y)}\right)
=
\max(a,b).
}
\]
清分母、齐次化并做和差回代即可得到整数多项式对实现两个输入 \(F\) 值的 min/max。

\section{批量最小值与第一极小位置}
固定 \(m\ge1\)，令
\[
H_k=\sum_{i=0}^{m-1}3^i x_i^k,\qquad k\ge m.
\]
写
\[
a_i=v_3(x_i),\qquad
a=\min_i a_i,
\]
并令
\[
i_0=\min\{i:a_i=a\}.
\]
第 \(i\) 项的赋值为
\[
ka_i+i.
\]
若 \(a_i>a\)，其相对 \(i_0\) 项的优势至少为
\[
k-(m-1)\ge1.
\]
若 \(a_i=a\) 而 \(i>i_0\)，编号偏移已经严格更大。

所以 \(i_0\) 项唯一最小：
\[
\boxed{v_3(H_k)=ka+i_0.}
\]
相邻幂次给出
\[
\boxed{
v_3(H_{m+1}/H_m)=a,
}
\]
以及
\[
\boxed{
v_3(H_m^{m+1}/H_{m+1}^m)=i_0.
}
\]

\section{有限整数热带表达式}
称整数热带多项式为
\[
T(\mathbf t)=\min_j(c_j+\mathbf a_j\cdot\mathbf t)
\]
其中 \(c_j\in\mathbb Z,\mathbf a_j\in\mathbb Z^m\)。

由整数常数、坐标、加减、min/max 有限复合得到的函数，恰好可以写成
\[
\boxed{g=T_1-T_2}
\]
其中 \(T_1,T_2\) 为整数热带多项式。

关键恒等式是：若
\[
f=A-B,\qquad g=C-D,
\]
则
\[
\min(f,g)
=
\min(A+D,C+B)-(B+D),
\]
而热带多项式的和仍是热带多项式。

\section{分类定理：构造方向}
在 \(x_i\) 坐标中：
\[
t_i+t_j\leftrightarrow x_ix_j,
\qquad
t_i-t_j\leftrightarrow x_i/x_j,
\]
\[
\min(t_i,t_j)\leftrightarrow\mu(x_i,x_j),
\]
\[
\max(t_i,t_j)\leftrightarrow x_ix_j/\mu(x_i,x_j),
\]
\[
r\leftrightarrow3^r.
\]
所有基本有理函数在 \((\mathbb Q^\times)^m\) 上无零点和极点。

有限复合得到固定有理函数 \(R_g\) 满足
\[
v_3(R_g(\mathbf x))=g(v_3(x_1),\ldots,v_3(x_m)).
\]
统一清分母、多重齐次化，并令
\[
P=L(\widetilde A+\widetilde B),\qquad
Q=L(\widetilde A-\widetilde B)
\]
即可回到整数多项式对，同时保持输出在定义域。

\section{有限单位网格与插值界}
设 \(C(Y_1,\ldots,Y_\ell)\in\mathbb Z[\mathbf Y]\) 非零，每个变量次数不超过 \(D\)。令
\[
S_D=\{1,4,7,\ldots,1+3D\}.
\]
记
\[
\tau(C)=\min_\alpha v_3(c_\alpha),
\]
\[
\eta(C)=\min_{\mathbf y\in S_D^\ell}v_3(C(\mathbf y)).
\]
则
\[
\boxed{
0\le
\eta(C)-\tau(C)
\le
\ell(D+v_3(D!)).
}
\]
下界来自超度量不等式。

一变量节点 \(y_j=1+3j\) 的 Lagrange 基分母为
\[
3^D\prod_{h\ne j}(j-h),
\]
其赋值至多
\[
D+v_3(D!).
\]
对 \(\ell\) 个变量张量化即可得到上界；否则所有网格值都过深会通过插值迫使全部系数比 \(\tau(C)\) 更深，矛盾。

\section{分类定理：必要性}
设固定整数多项式对 \((P,Q)\) 全域实现 \(g\)，并令
\[
A=P+Q,\qquad B=P-Q.
\]
把输入改写成 \(u_i=n_i+d_i,v_i=n_i-d_i\)，乘上固定的 \(2\) 的幂后可把 \(A,B\) 视为整数多项式。

固定
\[
\mathbf t=(t_1,\ldots,t_m)\in\mathbb Z^m
\]
并取
\[
u_i=2\cdot3^{t_i^+}s_i,\qquad
v_i=2\cdot3^{t_i^-}r_i,
\]
其中 \(s_i,r_i\in S_D\)。这给出实际整数输入并精确满足
\[
F(X_i)=t_i.
\]

对每个固定网格点 \(\omega\)，多项式展开中的每个单项赋值都是
\[
v_3(c_\alpha)+
\sum_i a_i t_i^+
+
\sum_i b_i t_i^-,
\]
属于有限整数热带表达式。

按这些有限仿射候选的比较关系划分 \(\mathbb Z^m\)。最低项唯一时赋值即为最低项；并列时提取共同幂，归一化后的单位系数和是固定整数。若完全抵消则删除这一组，若不完全抵消则只产生固定整数修正。项数有限，因此有限递归后得到：
\[
\mathbf t\mapsto v_3(A(\mathbf t;\omega))
\]
是有限整数仿射分区函数，也就是有限整数热带表达式。\(B\) 同理。

定义
\[
\eta_A(\mathbf t)=\min_{\omega}v_3(A(\mathbf t;\omega)),
\qquad
\eta_B(\mathbf t)=\min_{\omega}v_3(B(\mathbf t;\omega)).
\]
全域实现意味着对每个 \(\omega\)
\[
v_3(A(\mathbf t;\omega))
=
v_3(B(\mathbf t;\omega))+g(\mathbf t).
\]
所以
\[
\eta_A(\mathbf t)
=
\eta_B(\mathbf t)+g(\mathbf t),
\]
即
\[
\boxed{g=\eta_A-\eta_B.}
\]
有限最小值保持在有限整数热带表达式类中，必要性得证。

\section{Lipschitz 边界与条件开关}
整数仿射函数
\[
L(\mathbf t)=c+\mathbf a\cdot\mathbf t
\]
满足
\[
|L(\mathbf t)-L(\mathbf s)|
\le
\|\mathbf a\|_1
\|\mathbf t-\mathbf s\|_\infty.
\]
有限 min/max 保持最大斜率界，加减只会相加常数。因此每个可实现 \(g\) 都满足某个
\[
\boxed{
|g(\mathbf t)-g(\mathbf s)|
\le
L\|\mathbf t-\mathbf s\|_\infty.
}
\]

所以
\[
G(t,s)=
\begin{cases}
s,&t\ge0,\\
0,&t<0
\end{cases}
\]
不可实现，因为
\[
|G(0,M)-G(-1,M)|=M
\]
而输入距离恒为 \(1\)。

固定 \(B\ge0\) 时可以实现截断开关。令
\[
H(t)=\min(1,\max(0,t+1)),
\]
则
\[
H(t)=\mathbf1_{\{t\ge0\}}
\quad(t\in\mathbb Z).
\]
定义
\[
G_B(t,s)
=
\min(s_+,BH(t))
-
\min((-s)_+,BH(t)),
\]
便得到
\[
G_B(t,s)=
\begin{cases}
\max(-B,\min(s,B)),&t\ge0,\\
0,&t<0.
\end{cases}
\]

\section{一般素数 \(p\)}
对任意素数 \(p\)，
\[
\mu_p(x,y)=
\frac{x^3+py^3}{x^2+py^2}
\]
满足
\[
\boxed{
v_p(\mu_p(x,y))
=
\min(v_p(x),v_p(y)).
}
\]
证明与 \(p=3\) 完全相同：候选赋值为 \(3a,1+3b\) 和 \(2a,1+2b\)，永不相等。

因此这套 min/max 表达骨架不是 \(p=3\) 特有现象。

\section{文献边界与状态}
热带有理函数与分段线性函数的 min/max 表示已有独立文献。本文不声称一般热带表示理论新颖；当前固化的是 \(F_{3,D}\) 全域固定整数多项式模型中的精确算术实现等价。

本稿状态为 \texttt{PAPER-AUDITED}，尚未 Lean 形式化。

\begin{thebibliography}{9}
\bibitem{TW}
Ngoc M. Tran and Jidong Wang,
\emph{Minimal Representations of Tropical Rational Functions},
Algebraic Statistics 15 (2024), 27--59; arXiv:2205.05647.

\bibitem{KMPS}
Christoph Koutschan, Bernhard Moser, Anton Ponomarchuk and Josef Schicho,
\emph{Representing Piecewise Linear Functions by Functions with Small Arity},
Applicable Algebra in Engineering, Communication and Computing 36 (2025), 595--610; arXiv:2305.16933.
\end{thebibliography}
\end{document}
